\section{Objectives and signal representation}
This laboratory reviews discrete-time signals and systems through Scilab
functions for time shifting, folding, addition, pointwise multiplication,
and convolution. Each function is illustrated with an example, its source
code, and graphical or numerical output.

A finite sequence is represented by a vector \texttt{xn} and an integer
\texttt{xorigin}. If \texttt{xorigin} is within the stored vector, it is
the position of the sample at $n=0$. Since Scilab indices start at one,
the time corresponding to vector position $i$ is
\begin{equation}
  n_i=i-\texttt{xorigin},\qquad i=1,\ldots,N,
\end{equation}
where $N$ is the vector length. Thus the time axis is constructed as
\begin{lstlisting}[numbers=none]
nx = (1:length(xn)) - xorigin;
\end{lstlisting}
Samples outside the stored interval are assumed to be zero. After a time
shift, the origin parameter may fall outside $1,\ldots,N$; it still
defines the time coordinates and must not be used directly as a vector
index in that case.

For example, \texttt{xn = [1, -2, 3, 6]} and \texttt{xorigin = 3}
represent $x(-2)=1$, $x(-1)=-2$, $x(0)=3$, and $x(1)=6$.

The plots use vertical bars to represent discrete samples. Input and
output plots share the same axis limits, with a line thickness of three
pixels, so that their time positions can be compared directly.
